\documentclass[11pt,a4paper]{article}
\usepackage[margin=24mm]{geometry}
\usepackage{amsmath,amssymb,booktabs,longtable,array,hyperref}
\hypersetup{hidelinks}
\newcommand{\La}{\textbf{(a)}\ }
\newcommand{\Lc}{\textbf{(c)}\ }
\newcommand{\Ld}{\textbf{(d)}\ }
\newcommand{\scr}{\texttt{audit\_\allowbreak defq.py}}
\newcommand{\file}[1]{\texttt{\detokenize{#1}}}
\newcolumntype{L}[1]{>{\raggedright\arraybackslash}p{#1}}
\emergencystretch=3em
\title{Independent audit of defect charges and fusion:\\conventions, boundary traces and the scope of the obstructions}
\author{Quantum Boussinesq solitons project}
\date{4 November 2026}
\begin{document}
\maketitle
\begin{abstract}
\Lc The classical charges, continuum momentum trace, linear UV divergence and
KP growing-mode counts survive independent checks. The scattering parameter
is $s=\sigma_{\rm vac}$; the spectral parameter is legitimately
$c=\eta\sigma_{\rm vac}$. These are different quantities. A concrete extra
$\eta$ in the classical script turns its purported absorption test into
ordinary transmission. The printed energy-trace sign, negative-wavenumber
asymptotic transmission, threshold growth rate and pole argument against
root shifts are incorrect. Independent quadratures reproduce the numerical
canonical-charge holonomies and absorbed-momentum remainders.
\Ld Those examples and a conditional telescoping argument do not establish
a universal physical quantum-fusion no-go theorem.
\end{abstract}

\section{Scope, method and verdicts}
Labels retain the repository meanings: (a) established in literature,
(b) literature heuristic, (c) project derivation or computational evidence,
(d) unestablished physical identification or research judgment.
\Lc This audit reads the package source/PDF and three scripts, the relevant
parts of \file{paper/review_draft1.tex}, the charges/action and frames/defects
audits, \file{audit_phase2i.tex}, and the CZ repository notes. All independent
checks are in \file{phase2/scripts/audit_defq.py}. No package checker is
imported or used to supply numerical data to this checker. The original
three scripts were separately run to test their exit status and success messages.

\Lc The independent methods are action variation with a free vacuum
constant, integration of the momentum-flux one-form, polynomial reduction
on a KP edge, logistic-coordinate spatial quadrature, fresh linear sewing
solves, direct determinants away from their zeros, and rational kernel
quadrature followed by numerical differentiation. SymPy checks are exact;
mpmath uses 70 decimal digits, normally with absolute tolerance $10^{-55}$.
Finite-cutoff tests retain their stated asymptotic error.

\begin{longtable}{@{}L{5.1cm}L{2.3cm}L{7.9cm}@{}}
\toprule Claim & Verdict & Independent evidence and limitation\\\midrule\endhead
Scattering root and growth parameter & Qualified (c) & Correct distinction in most formulas; extra $\eta$ corrupts the absorption input.\\
Four classical charges and sewing fluxes & Confirmed (c) & Action variation, integrated flux one-form, both orientations; bulk equations and sewing required.\\
KP-edge charge factorisation & Confirmed (c) & Exact polynomial quotient and both orientations.\\
Package's numerical absorption test & Refuted (c) & Transmission factor $1.58257569495584$, not zero. Independent genuine absorption tests pass.\\
Corrected determinant and branch switch & Confirmed (c) & Full action-derived matrices for both orientations.\\
Growing counts; no stable--stable edge & Confirmed (c) & $(0,1,2)$ generically, $(0,1,1)$ at switch; printed threshold growth value refuted.\\
Quadratic momentum trace & Confirmed (c) & Fresh sewing modes, polarization and integral. Complete physical identification qualified (d).\\
Printed energy trace and asymptotic $z$ & Refuted (c) & Energy trace has a plus sign; signed-wavenumber limits coincide.\\
Linear Casimir divergence; boundary term & Confirmed (c) & Independent cutoff quadrature and UV slope; finite scheme not selected.\\
Only parameter shifts allowed as counterterms & Qualified & Restricted potential test passes (c); exhaustive quantum classification absent (d).\\
Canonical-charge holonomy no-go & Qualified & Examples reproduced (c); universal nonvanishing unproved. Telescoping valid at fixed edge data.\\
Absorbed-soliton momentum remainder & Qualified & Identity and nonzero examples confirmed (c); physical frame conditional (d).\\
Pole proof excluding root shifts & Refuted (c) & Alleged simple poles have zero residue on the soliton root locus.\\
Complete physical one-loop factorisation & Qualified (d) & Continuum identities do not construct the quantum boundary theory.\\
Every original checker prints success & Confirmed (c) & All three exit zero and print \file{ALL CHECKS PASSED}, despite the discrepancies.\\\bottomrule
\end{longtable}

\section{Two parameters, and one actual convention error}
\La The source defines $\sigma_0=sc$ in its dimensional notation; with bulk
wave speed one, its scattering parameter is $s=\sigma_0$ \cite[Sec.~5.1]{CZ}.
\Lc To prevent collision with that notation, write
\begin{equation}
 s=\sigma_{\rm vac},\qquad q_D=-\frac{s}{2\sqrt3},\qquad
 c=\eta\sigma_{\rm vac},\qquad \eta=\pm1.
 \label{eq:dictionary}
\end{equation}
Here $c$ denotes only the spectral combination, not the dimensional wave
speed. Fresh action-derived sewing solves give continuum transmission
\begin{equation}
 z(q;s,\eta)=\left(
 \frac{\operatorname{sgn}(q)\sqrt{1+q^2}+s-i\sqrt3q}
 {\operatorname{sgn}(q)\sqrt{1+q^2}+s+i\sqrt3q}\right)^\eta.
 \label{eq:phonon}
\end{equation}
\Lc The checker solves one sewing equation for $z$, checks all four residuals,
and only then compares with \eqref{eq:phonon}. The decaying-mode determinant
instead depends on $c=\eta s$.

\Lc The package's classical symbolic action uses bare $s=\sigma_0$ correctly.
Its determinant test is at $\eta=+1$, where the two parameters coincide.
Its fusion checker correctly uses $s=-2\sqrt3q_D$ in dipole charges and
$c=-2\sqrt3\eta q_D$ in growing-mode counts. Replacing $c$ by $s$ in the
spectrum would therefore introduce an error. The exception is the classical
checker's absorption input:
\begin{quote}\small
\file{s_abs = -2 * r3 * q1 * eta_abs}
\end{quote}
\Lc The receiving function uses \file{qD = -s_val / (2 * r3)}, treating
\file{s_val} as the bare parameter. Its actual inputs give
\begin{equation}
 q_1=\frac{\sqrt3}{4},\quad
 q_2=\frac{-q_1+\sqrt{1-3q_1^2}}2,\quad \eta=-1,
 \qquad s_{\rm passed}=+1.5,\quad s_{\rm required}=-1.5.
\end{equation}
The resulting $z=1.58257569495584001$ is regular transmission. The reported
conservation error is real, but does not verify absorption: the original
script never asserts $z=0$. Our checker tests genuine $z=0$ solutions for
both orientations and their four conserved charges.

\Lc A second diagnostic reproduces the review's example:
$\alpha=0.4$, $\epsilon=1$, $\sigma_0=2.4$, $\eta=-1$ gives
$z=1.5097108339861$. Inserting $\eta\sigma_0$ in the scattering root instead
gives $0.37356503047145$; the maximum sewing residual is
$1.1104178917082$ in project potential units (three times this in CZ field
units). The correct solution passes the $10^{-55}$ test.

\Lc The correctly implemented charges, growth rates, triangle counts and
algebraic obstructions need no orientation correction. The numerical
absorption example does. A literal substitution of
$-\eta\sigma_0/(2\sqrt3)$ for the root at fixed fields changes transmission
and root-labelled charge assignments. It is not a harmless renaming.

\section{Classical charges rederived from the action}
\La The starting action and sewing variation are CZ's type-I construction
\cite[Secs.~2--3]{CZ}. \Lc Put $a=\varphi_L$, $b=\varphi_R$, $x=a_x$,
$y=b_x$, and $A=\eta/\sqrt3$. The explicit dictionary is
$U=-3a$, $V=-3b$, $\mathcal L_{CZ}=9\mathcal L$, $D=D_{CZ}/9$.
The independently integrated functions satisfy
\begin{align}
 \sigma&=s-(a-b)/A,&
 \rho(\sigma)&=\rho_0+\frac{\sigma-\sigma^3/3}{2A},\\
 \tau'(\sigma)&=-\sigma\rho(\sigma),&
 h'(\sigma)&=-\rho(\sigma),\qquad \tau(s)=h(s)=0.
\end{align}
The vacuum constraint is $\rho_0=-(s-s^3/3)/(2A)$. The gauge
$f=g=-\sigma^2/(4A)$ satisfies $f_V-g_U=\sigma$; the factor-two discrepancy
in the source's example gauge is avoided by enforcing that constraint.
\Lc Conversion of the action gives
\begin{align}
 D&=a_tF+b_tG-E_D,\\
 F&=A(x+2y)+\frac{\sigma^2}{12A},&
 G&=A(2x+y)+\frac{\sigma^2}{12A},\\
 E_D&=-A\sigma(x^2+xy+y^2)+\frac{\rho}{3}(x+y)-\frac{\tau}{9},\\
 P_D&=-\frac A2(x^2+4xy+y^2)+\frac h9,\\
 \Pi_D&=F+G=3A(x+y)+\frac{\sigma^2}{6A},&
 M_D&=b-a=A(\sigma-s).
 \label{eq:charges}
\end{align}
\Lc The checker first recovers $P_D$ by integrating the one-form required
by momentum flux, then compares with \eqref{eq:charges}. This checks that
CZ's intermediate $-U_xF_{CZ}-V_xG_{CZ}$ terms must not be subtracted again.
Variation gives
\begin{align}
 a_{xx}&=D_x,&a_{xxx}&=x-3x^2-D_a+\dot F,\\
 b_{xx}&=-D_y,&b_{xxx}&=y-3y^2+D_b-\dot G,
\end{align}
where $D_x,D_y$ are partial derivatives with respect to boundary jet variables.
\Lc The bulk flux primitives used in the independent check are
\begin{align}
 B_E&=\varphi_t(\varphi_x-3\varphi_x^2-\varphi_{xxx})+\varphi_{xt}\varphi_{xx},\\
 B_P&=-\tfrac12\varphi_t^2-\tfrac12\varphi_x^2+2\varphi_x^3
            +\varphi_x\varphi_{xxx}-\tfrac12\varphi_{xx}^2,\\
 B_\Pi&=\varphi_x-3\varphi_x^2-\varphi_{xxx},\qquad B_M=\varphi_t.
\end{align}
\Lc Symbolically, $B_Q(L)-B_Q(R)+\dot Q_D=0$ for all four charges after
sewing substitution, for arbitrary remaining jets and both orientations.
``Off-shell'' means no special soliton solution was imposed on the boundary
jets. Conservation of bulk integrals still uses the bulk equation; it is
not an identity for unrestricted off-shell bulk fields. Fixed asymptotic
potentials are required for the mass statement.

\Lc For an edge $p\to q$, reduction modulo $p^2+pq+q^2-1/4$ gives the
four charge jumps with signed $k=p-q$ and $v=\sqrt3(p+q)$. Restoring positive
$k=|p-q|$ yields exactly
\begin{equation}
 \Delta(E_D,P_D,\Pi_D,M_D)
 =\eta\operatorname{sgn}(p-q)(E,P,\Pi,M)_{\{p,q\}}.
\end{equation}
\Lc The time-dependent check uses $h=(1+e^{-k(x-vt)})^{-1}$,
$\varphi=2kh$, integrating bulk densities in $h$ separately on the two
half-lines. All four totals agree at $t=-3,0.27,4$ to $10^{-55}$,
for transmission and actual absorption at both orientations.

\section{Quadratic traces: the resolved question and its limits}
\Lc Taylor expansion of the action gives the package's quadratic operators,
including $-\sigma_0(a-b)/2$ in $F_1,G_1$. The checker contracts them with
fresh complex sewing amplitudes by quarter-cycle polarization, using the
vacuum covariance $\langle a^2\rangle_q=1/(4\pi\omega)$ and
$\omega=|q|\sqrt{1+q^2}$. For $q>0$,
\begin{align}
 \langle P_{D,2}\rangle_q+\langle P_{D,2}\rangle_{-q}
 &=\frac{q}{4\pi}\{\theta'(q)-\theta'(-q)\}\\
 &=\frac{\sqrt3\eta q\{2(\sigma_0^2-2)q^2+\sigma_0^2-1\}}
 {\pi\sqrt{1+q^2}\{16q^4+4(\sigma_0^2+2)q^2+(\sigma_0^2-1)^2\}}.
 \label{eq:ptrace}
\end{align}
\Lc The maximum residual over five vacuum parameters, three positive
wavenumbers and both orientations is $1.05\,10^{-71}$. Independent integration
of this rational trace agrees with dipole momentum from compact-interval
kernel quadrature to $10^{-55}$ at two stable oriented vacua.

\Lc This resolves the earlier missing \emph{continuum quadratic operator}
calculation. The charges/action audit asked for this calculation explicitly.
Here an action-derived operator is compared with independently solved scattering data. It is not a formula
compared with itself. \Ld The complete physical identification still needs
a boundary domain, canonical mode completeness, discrete/threshold sectors,
a renormalised vacuum and quantum Ward identities. Infinite-volume plane-wave
bulk subtraction alone does not settle finite-volume normalisation.
On unstable vacua these are formal continuum traces, not constructed
quantum ground-state expectations.

\paragraph{Printed discrepancies.}
\Lc Let $K=a_tF_1+b_tG_1$. Direct contraction gives
\begin{equation}
 \langle E_{D,2}+J_{\rm bulk,2}\rangle_q=\langle K\rangle_q,
 \qquad \langle K\rangle_q+\langle K\rangle_{-q}
 =\frac{\omega}{4\pi}\{\theta'(q)+\theta'(-q)\}.
\end{equation}
\Lc For these plane waves $\langle J_{\rm bulk,2}\rangle=0$.
The minus sign in the package's boxed energy identity is false. Its code
asserts \file{2 * symp_bdry_sym - dth_even_E == 0}, where the former term
is $\langle K\rangle_{\rm sym}$, but prints the opposite sign as a pass.
The audit asserts the nonzero discrepancy with that printed identity.

\Lc Also, \eqref{eq:phonon} gives
\begin{equation}
 z(+\infty)=z(-\infty)=e^{-2\pi i\eta/3},
\end{equation}
not $e^{-2\pi i\eta\operatorname{sgn}(q)/3}$. This does not change phase
derivatives, but matters to the cutoff endpoint gauge. Finally, the printed
real mode of amplitude $1/\sqrt{\pi\omega}$ has ordinary time average
$1/(2\pi\omega)$, not $1/(4\pi\omega)$. The latter is the standard vacuum
contraction actually used by the formulas; the text must specify an extra
half weight or use amplitude $1/\sqrt{2\pi\omega}$. Our checker fixes the
covariance explicitly.

\section{Spectrum and obstruction (i)}
\Lc Fresh differentiation of the nonlinear action and substitution of
$e^{gt\pm\kappa x}$ gives, in the package's row ordering,
\begin{align}
 \frac{\det M}{(\kappa_1-\kappa_2)^2}
 &=-\frac23 S^2(S^2-1)P_-(S;c)P_+(S;c),\\
 P_\mp(S;c)&=S^2-\frac{\sqrt3\mp1}{2}cS+\frac{c^2-2\mp\sqrt3}{2},\\
 c&=\eta\sigma_0,\qquad S=\kappa_1+\kappa_2,\qquad g=(S^2-1)/2.
\end{align}
\Lc The checker solves the bulk quartic for decaying exponents and compares
the unreduced determinant at sixteen points away from its zeros. Each
candidate admissible rate is also checked in the full matrix. This confirms
\file{audit_frames_defects.tex} and \file{audit_phase2i.tex}, including the
extension to $\eta=-1$.

\Lc Selecting $S>1$ gives zero growing modes for $c<-1$, one for
$-1<c<1$, two for $1<c<2$ except at $c=\sqrt3$, where one is marginal.
The switch from $P_-$ to the second $P_+$ root is correct, but
\begin{equation}
 g_{\rm remaining}(\sqrt3)=\frac{\sqrt3}{4}
 =0.4330127018922193\ldots\ne\frac{2\sqrt3-1}{8}.
\end{equation}
\Lc The last expression is the incorrect printed value, not a change to the
count. For $\eta=-1$, $\sigma_0=-1.9$, the actual rates are
$0.2002565225612746$ and $0.4968939274851985$; using bare $s=-1.9$
instead of correct $c=+1.9$ would wrongly predict zero.

\Lc Three distinct real KP roots give one spectral parameter in each of
$(-2,-1)$, $(-1,1)$ and $(1,2)$, as follows from their common cubic's extrema.
Thus the sorted counts are $(0,1,2)$ generically and $(0,1,1)$ at the branch
switch. The independent checker covers 99 nondegenerate angles and both
orientations, in addition to the determinant tests. No edge connects two
vertices without a growing mode. Marginal points and repeated-root triples
need separate treatment; these checks are not a proof of nonlinear stability
or full spectral completeness. \Ld The usual restriction to stable
asymptotic defect states does not rule out a separately constructed resonance
continuation.

\section{UV divergence and boundary counterterms}
\Lc Set the phase of \eqref{eq:phonon} to zero at each infinity and put
$T=\sqrt{1+\Lambda^2}$, $S_D=-\eta\sqrt3\sigma_0/2$. Independent phase
differentiation and quadrature confirm
\begin{align}
 E_{\rm dens}&=\frac{S_D}{2\pi}(T-1)+\Delta E_D^{\rm ren}
                                  -\frac{S_D}{2\pi}+O(\Lambda^{-1}),\\
 B_D&=-\frac{S_D}{2\pi}T+O(\Lambda^{-1}),\\
 E_{\rm mode}=E_{\rm dens}-B_D
 &=\frac{S_D}{\pi}(T-1)+\Delta E_D^{\rm ren}+O(\Lambda^{-1}).
\end{align}
\Lc The independent cutoffs are $10^3,10^5$; the finite term is checked by a
separate compact-interval kernel quadrature. The limiting density slope is
also checked at $q=10^{20}$. At either stable vacuum
$(\eta,\sigma_0)=(1,-1.5)$ or $(-1,1.5)$ the limiting density and mode slopes
are $0.206748335783\ldots$ and $0.413496671566\ldots$.
The original renormalisation script's header claims a different slope and
polynomial boundary-operator divergences; those stale comments contradict
its actual calculations.

\Lc The open-half-line bulk tadpole vanishes on constant vacua and cannot
cancel this divergence. The proposed boundary term decomposes as
\begin{equation}
 -\frac{S_D}{\pi}(T-1)
 =\frac{\eta\sqrt3 s}{2\pi}(T-1)+\frac3{2\pi}(T-1)M_D.
\end{equation}
\Lc This confirms the boundary-jump completion of the total derivative,
up to a vacuum-independent reference term. It establishes a boundary
subtraction requirement in this prescription without selecting a finite part.

\Lc Perturbing the action by a potential with unchanged derivative couplings
changes momentum flux by $-(x-y)\delta\tau'$ in the corresponding units.
Independent jets cannot match $(a_t-b_t)\delta h'$ unless both derivatives
vanish. The checker derives this change from perturbed sewing equations.
\Ld This restricted result is not an exhaustive classification of quantum
boundary operators, derivative counterterms or anomalous conservation laws.
Parameter shifts within the CZ family are available; their claimed necessity
for every integrable quantum renormalisation is not proved. A fully
established physical one-loop transmission theory does not follow.

\section{Obstruction (ii): holonomy and the no-go theorem's scope}
\Lc The independent charge calculation integrates rational phase kernels,
not the package's logarithmic charge functions. With
$\alpha=\arccos|v|$, $\epsilon=\operatorname{sgn}v$,
$\phi_j=\alpha+(2\pi/3,-2\pi/3,0)_j$, $w_j=(1,1,-2)_j$, it uses
\begin{align}
 I_j&=\int_0^1\frac{dz}{1-\cos^2\phi_j\,z^2},\\
 E_m&=\frac1{2\pi}\sum_j w_j\sin\phi_j(1+\cos^2\phi_j I_j),&
 P_m&=\frac\epsilon{2\pi}\sum_j w_j\sin\phi_j\cos\phi_j I_j.
\end{align}
\Lc Analogous dipole quadratures use
$\psi_\pm=\pm2\pi/3-\arccos(\sigma_0/2)$, weights $(-\eta,+\eta)$.
They reproduce edge differences and zero mode-charge holonomy to $10^{-55}$.
Canonical charges are obtained independently by differentiation:
\begin{equation}
 L(v)=vP_m(v)-E_m(v),\qquad P_c=L'(v),\qquad E_c=vP_c-L.
\end{equation}
\Lc The three published examples agree:
\begin{center}
\begin{tabular}{rrr}\toprule
Seed root & $\sum\rho_iE_c$ & $\sum\rho_iP_c$\\\midrule
$-0.20$ & $0.3175373286355155$ & $0.1655848943300202$\\
$0.10$ & $0.4898832009097553$ & $-0.1111527059147391$\\
$0.45$ & $-0.5020709673515473$ & $0.0996575496842790$\\\bottomrule
\end{tabular}
\end{center}
\Lc Two additional triangles also have nonzero vectors. Reversing $\eta$
reverses each directed holonomy. The symmetric nondegenerate triple
$\{0,1/2,-1/2\}$ has holonomy $(0.5445577069714241,0)$: individual component
inequalities need care even when the vector obstruction persists.

\Lc The original script tests only three seeds for canonical nonvanishing.
Its 594-edge scan concerns stability, not holonomy. No analytic proof of
nonvanishing for every generic triangle is provided. Our additional samples
are also not such a proof. Examples are confirmed; the universal premise is
qualified.

\Lc The counterterm argument is rigorous \emph{conditionally}: for any
single-valued vertex function $C$,
\begin{equation}
 [C(q_2)-C(q_1)]+[C(q_3)-C(q_2)]+[C(q_1)-C(q_3)]=0.
\end{equation}
\Lc At a triangle with demonstrated nonzero fixed canonical edge vector,
no vertex counterterm balances all three edges. This also holds for divergent
counterterms at a common finite regulator. It does not cover changes to
absorption roots, rapidities, the charge dictionary or the state space.
The package retains a kinematic escape: since
$E_c-E_m=v(P_c-P_m)$ and $E_{\rm cl}'=vP_{\rm cl}'$,
$\delta v=(P_m-P_c)/(dP_{\rm cl}/dv)$ matches both first-order charges
when its denominator is nonzero. The audit verifies these derivatives.
\Ld Quantum sewing must still derive such a shift and identify the physical
charges. ``Regardless of boundary counterterms'' is justified only while
keeping the fixed-edge, single-valued-state assumptions.

\section{Obstruction (iii) and the failed pole argument}
\Lc Independent Hirota-polynomial and root-ratio evaluations confirm
summation by parts. Write $L_i=\log|z_t(3;q_i,\eta)|$,
$w_i=\Delta P(S_i)/k_i$. Then
\begin{equation}
 \sum_i\rho_{{\rm tree},i}\Delta X_i\Delta P(S_i)
 =\sum_i w_i(L_{i+1}-L_i)=\sum_i(w_{i-1}-w_i)L_i.
\end{equation}
\Lc At the package's test point the remainders are
$-0.7742169719045726$ (mode momentum) and
$-1.737010868355514$ (canonical momentum). The checker tests both
orientations and both test-soliton directions, as well as vanishing of the
test-soliton shift sum. The absorbed momentum term is a real algebraic issue
in that comparison. \Ld Its unavoidable physical interpretation still
requires the charge/frame dictionary; nonconstant weights do not prove
nonzero values for every possible test soliton.

\Lc The subsequent pole argument excluding root shifts fails. Its proposed
simple pole is proportional to $(\ell_3-\ell(q))/(p_{3,a/b}-q)$, but
$\ell(p)=3\sqrt3(4p^3-p)$ and $\ell_3=\ell(p)$ at either soliton root.
Exactly,
\begin{equation}
 \frac{\ell(p)-\ell(q)}{p-q}
 =3\sqrt3\{4(p^2+pq+q^2)-1\}.
\end{equation}
\Lc This is a polynomial with zero pole residue. Differentiating the
root-dependent coefficient of the phase still gives logarithms, so the
asserted contrast with the logarithmic remainder is false. The checker
asserts this cancellation and hence the discrepancy.
\Ld This refutes the stated proof; it does not prove that consistent
root shifts exist. That functional cancellation and quantum sewing problem
remain unsolved here.

\section{Circularity, execution and failures recorded}
\Lc The original nonlinear flux identities, determinant comparison away
from zeros, and momentum-trace comparison are meaningful tests of distinct
expressions. The latter inserts a known transmission factor and proves its
symbolic identity only at $\eta=+1$; our fresh action-derived modes close
that gap. The absorbed-momentum comparison uses distinct Hirota and root
expressions. Neither is a literal self-comparison. Conversely, holonomy
uses assumed charge functions, not a field effective-action or Ward-identity
calculation; it cannot independently establish their physical meaning.

\Lc Important unsupported passes are the energy trace printed with the
opposite sign from its assertion, absorption without a zero-transmission
assertion, an untested pole argument, and physical factorisation without a
quantum-sewing computation. Stale comments contradict final calculations.
All three original scripts nevertheless exit zero and print
\file{ALL CHECKS PASSED}.

\paragraph{Failures recorded.}
\begin{enumerate}
\item \Lc Refuted package statements: purported numerical absorption,
energy-trace sign, negative-$q$ asymptotic transmission, threshold growth
value, and the simple-pole exclusion of root shifts. The real-mode amplitude
also needs an explicit half-weight convention. Audit assertions encode the
discrepancies, not the false equalities.
\item \Lc The first independent UV run failed its mode-slope assertion
because it used the printed sign-dependent transmission limit as phase
reference. Fresh sewing modes have the same constant at both infinities.
Correcting that reference removed the spurious cutoff contribution without
changing the target slope or tolerance.
\item \Lc A symbolic boundary-counterterm check failed because $\eta^2=1$
was substituted before expansion exposed the square. Expanding first fixed
the reduction without changing the identity or tolerance.
\item \Lc An initial search used an invalid regular expression; a literal
search fixed it. The first report-assembly command used a repository-relative
path while already inside \file{phase2/}, so no file was written; the corrected
path fixed assembly. These were tool/path errors, not mathematical evidence.
The first PDF build also exposed an overfull table and an overfull filename
line; column sizing and prose wrapping were corrected. One render attempt
started before the PDF writer had finished and reported an incomplete cross-reference
table; waiting for build completion and rerendering fixed it. No other independent
mathematical assertion failed during development.
\end{enumerate}

\Lc The final independent checker prints \file{ALL CHECKS PASSED} after
its exact identities, high-precision checks and refutation assertions.
Only the three requested \file{audit_defq} deliverables are retained in the
repository. No existing file or git state was changed. The full repository
suite was not rerun during this targeted audit with concurrent agents.
The PDF was built by two \file{pdflatex} passes from \file{phase2/} and
rendered for visual inspection.

\paragraph{Reading scope and references.}
\La The source below supplies the CZ action, variational sewing and
scattering-parameter definition. Relevant full-text sections 2, 3 and 5.1
were read online for this audit, not just the abstract; the entire article
was not read. \Lc New formulas and corrections above have independent
script checks. No new FK attribution is made: the equal-mode cutoff is
audited as the package's prescription, without claiming an independent
reading of that literature.

\begin{thebibliography}{9}
\bibitem{CZ} E.~Corrigan and C.~Zambon,
\emph{Adding integrable defects to the Boussinesq equation},
\href{https://arxiv.org/html/2306.03772v1}{arXiv:2306.03772v1}.
Relevant full-text reading scope specified above.
\end{thebibliography}
\end{document}
